\documentclass[12pt, a4paper]{ctexart}
\usepackage{amsmath, amssymb, amsthm}
\usepackage{fancyhdr}
\usepackage{geometry}
\usepackage{enumitem}
\usepackage{graphicx}
\usepackage{tikz}
\usepackage{pgfplots}
\pgfplotsset{compat=1.18}

\geometry{left=2cm, right=2cm, top=2.5cm, bottom=2.5cm}
\setCJKmainfont{Hiragino Sans GB}
\setCJKsansfont{Heiti SC}

\pagestyle{fancy}
\fancyhf{}
\fancyhead[L]{培优专题·极点极线与射影几何}
\fancyhead[R]{测试卷}
\fancyfoot[C]{\thepage}
\renewcommand{\headrulewidth}{0.4pt}

\newcommand{\score}[1]{\hfill\makebox[2cm][l]{(#1分)}}

\begin{document}

\begin{center}
\Large\textbf{培优专题测试卷·极点极线与射影几何}

\normalsize
考试时间：120分钟\quad 满分：150分
\end{center}

\vspace{0.5cm}

\noindent\textbf{一、单选题（本大题共8小题，每小题5分，共40分。在每小题给出的四个选项中，只有一项是符合题目要求的。）}

\vspace{0.3cm}

\begin{enumerate}
\item 已知椭圆 $\dfrac{x^2}{4} + \dfrac{y^2}{3} = 1$，点 $P(1, 0)$ 关于此椭圆的极线方程为\score{5}
\begin{enumerate}[label=\Alph*.]
\item $x = 4$
\item $x = 2$
\item $x = 1$
\item $x = \dfrac{3}{2}$
\end{enumerate}

\item 设椭圆 $\dfrac{x^2}{a^2} + \dfrac{y^2}{b^2} = 1$（$a > b > 0$），点 $P(x_0, y_0)$ 在椭圆外，则 $P$ 关于椭圆的极线与椭圆的位置关系为\score{5}
\begin{enumerate}[label=\Alph*.]
\item 相交
\item 相切
\item 相离
\item 不确定
\end{enumerate}

\item 已知抛物线 $y^2 = 4x$，点 $F(1, 0)$ 为焦点，则 $F$ 关于此抛物线的极线为\score{5}
\begin{enumerate}[label=\Alph*.]
\item $x = -1$
\item $x = 1$
\item $y = -1$
\item $y = 1$
\end{enumerate}

\item 点 $P(3, 0)$ 关于椭圆 $\dfrac{x^2}{9} + y^2 = 1$ 的极线方程为\score{5}
\begin{enumerate}[label=\Alph*.]
\item $x = 1$
\item $x = 3$
\item $x = \dfrac{1}{3}$
\item $x = 9$
\end{enumerate}

\item 已知双曲线 $\dfrac{x^2}{4} - \dfrac{y^2}{5} = 1$，点 $P(4, 0)$ 关于此双曲线的极线方程为\score{5}
\begin{enumerate}[label=\Alph*.]
\item $x = 1$
\item $x = 2$
\item $x = \dfrac{1}{2}$
\item $x = 4$
\end{enumerate}

\item 设 $A, B$ 是椭圆 $\dfrac{x^2}{a^2} + \dfrac{y^2}{b^2} = 1$ 上两点，$P$ 是 $AB$ 的极点，若 $AB$ 过焦点 $F(c, 0)$，则 $P$ 位于\score{5}
\begin{enumerate}[label=\Alph*.]
\item 对应准线上
\item 长轴上
\item 短轴上
\item 另一焦点处
\end{enumerate}

\item 已知圆 $x^2 + y^2 = r^2$，点 $P(a, b)$ 在圆外，则 $P$ 关于圆的极线被圆截得的弦长为\score{5}
\begin{enumerate}[label=\Alph*.]
\item $\dfrac{2r\sqrt{a^2 + b^2 - r^2}}{a^2 + b^2}$
\item $\dfrac{2r^2}{\sqrt{a^2 + b^2}}$
\item $2\sqrt{a^2 + b^2 - r^2}$
\item $\dfrac{2r}{\sqrt{a^2 + b^2 - r^2}}$
\end{enumerate}

\item 设椭圆 $\dfrac{x^2}{4} + y^2 = 1$，$A, B$ 在椭圆上，$M$ 为 $AB$ 中点，$O$ 为原点。若 $k_{OA} \cdot k_{OB} = -\dfrac{1}{4}$，则 $M$ 的轨迹为\score{5}
\begin{enumerate}[label=\Alph*.]
\item 椭圆 $\dfrac{x^2}{2} + y^2 = \dfrac{1}{2}$
\item 椭圆 $x^2 + 2y^2 = 1$
\item 椭圆 $\dfrac{x^2}{2} + 2y^2 = 1$
\item 圆 $x^2 + y^2 = \dfrac{1}{2}$
\end{enumerate}

\end{enumerate}

\vspace{0.5cm}

\noindent\textbf{二、多选题（本大题共4小题，每小题5分，共20分。在每小题给出的选项中，有多项符合题目要求。全部选对得5分，部分选对得2分，有选错的得0分。）}

\vspace{0.3cm}

\begin{enumerate}[resume]
\item 关于极点极线，下列说法正确的有\score{5}
\begin{enumerate}[label=\Alph*.]
\item 若点 $P$ 在圆外，则 $P$ 的极线与圆相交
\item 若点 $P$ 在椭圆上，则 $P$ 的极线即为椭圆在 $P$ 处的切线
\item 若点 $P$ 在圆内（非圆心），则 $P$ 的极线在圆外
\item 极点与极线的关系是对偶的
\end{enumerate}

\item 已知椭圆 $\dfrac{x^2}{4} + \dfrac{y^2}{3} = 1$，$P(4, 0)$，过 $P$ 作两条切线切椭圆于 $A, B$ 两点，则\score{5}
\begin{enumerate}[label=\Alph*.]
\item $AB$ 的方程为 $x = 1$
\item $\triangle PAB$ 的面积为 $\dfrac{3\sqrt{3}}{2}$
\item $|PA| = |PB|$
\item $PA \perp PB$
\end{enumerate}

\item 关于圆锥曲线的极线，下列正确的有\score{5}
\begin{enumerate}[label=\Alph*.]
\item 椭圆 $\dfrac{x^2}{a^2} + \dfrac{y^2}{b^2} = 1$，点 $P(x_0, y_0)$ 的极线为 $\dfrac{x_0 x}{a^2} + \dfrac{y_0 y}{b^2} = 1$
\item 双曲线 $\dfrac{x^2}{a^2} - \dfrac{y^2}{b^2} = 1$，点 $P(x_0, y_0)$ 的极线为 $\dfrac{x_0 x}{a^2} - \dfrac{y_0 y}{b^2} = 1$
\item 抛物线 $y^2 = 2px$，点 $P(x_0, y_0)$ 的极线为 $y_0 y = p(x + x_0)$
\item 圆 $x^2 + y^2 = r^2$，点 $P(x_0, y_0)$ 的极线为 $x_0 x + y_0 y = r^2$
\end{enumerate}

\item 设椭圆 $\dfrac{x^2}{4} + y^2 = 1$，$A, B$ 为椭圆上两点，$P$ 为 $AB$ 的极点，下列说法正确的有\score{5}
\begin{enumerate}[label=\Alph*.]
\item 若 $AB$ 过原点，则 $P$ 在无穷远处
\item 若 $AB$ 平行于 $x$ 轴，则 $P$ 在 $y$ 轴上
\item 若 $AB$ 是椭圆的切线，切点为 $T$，则 $P = T$
\item $P$ 到椭圆中心的距离与 $AB$ 到中心的距离之积为 $a^2 b^2$（在适当的定向意义下）
\end{enumerate}

\end{enumerate}

\vspace{0.5cm}

\noindent\textbf{三、填空题（本大题共4小题，每小题5分，共20分。）}

\vspace{0.3cm}

\begin{enumerate}[resume]
\item 椭圆 $\dfrac{x^2}{9} + \dfrac{y^2}{4} = 1$，点 $P(3, 2)$ 关于此椭圆的极线方程为 \underline{\hspace{3cm}}。\score{5}

\item 抛物线 $y^2 = 8x$，点 $P(2, 4)$ 关于此抛物线的极线方程为 \underline{\hspace{3cm}}。\score{5}

\item 已知椭圆 $\dfrac{x^2}{4} + y^2 = 1$，点 $P(0, 2)$，过 $P$ 作椭圆的两条切线，切点分别为 $A, B$，则直线 $AB$ 的方程为 \underline{\hspace{3cm}}。\score{5}

\item 设 $A(x_1, y_1)$，$B(x_2, y_2)$ 是椭圆 $\dfrac{x^2}{a^2} + \dfrac{y^2}{b^2} = 1$ 上两点，$AB$ 的中点为 $M(x_0, y_0)$，则 $k_{AB} \cdot k_{OM} = $ \underline{\hspace{2cm}}（其中 $k_{OM}$ 为 $OM$ 的斜率）。\score{5}

\end{enumerate}

\vspace{0.5cm}

\noindent\textbf{四、解答题（本大题共6小题，共70分。解答应写出文字说明、证明过程或演算步骤。）}

\vspace{0.3cm}

\begin{enumerate}[resume]
\item （10分）已知椭圆 $\dfrac{x^2}{4} + y^2 = 1$，点 $P(4, 0)$。
\begin{enumerate}[label=(\arabic*)]
\item 求 $P$ 关于椭圆的极线方程；
\item 过 $P$ 作两条切线切椭圆于 $A, B$，求 $|AB|$ 及 $\triangle PAB$ 的面积。
\end{enumerate}

\item （12分）已知椭圆 $\dfrac{x^2}{6} + \dfrac{y^2}{2} = 1$，$A, B$ 为椭圆上两点，$M$ 为 $AB$ 中点。
\begin{enumerate}[label=(\arabic*)]
\item 若 $M(1, \dfrac{1}{2})$，求直线 $AB$ 的方程；
\item 证明：$k_{AB} \cdot k_{OM} = -\dfrac{1}{3}$；
\item 设 $P$ 为 $AB$ 关于椭圆的极点，求 $P$ 的轨迹方程。
\end{enumerate}

\item （12分）已知抛物线 $y^2 = 4x$，$A, B$ 为抛物线上两点，$F$ 为焦点。
\begin{enumerate}[label=(\arabic*)]
\item 证明：若 $AB$ 过 $F$，则 $AB$ 的极点在准线上；
\item 若 $A(t_1^2, 2t_1)$，$B(t_2^2, 2t_2)$，求 $AB$ 的极点坐标；
\item 若 $AB$ 的极点为 $P(-1, 2)$，求直线 $AB$ 的方程。
\end{enumerate}

\item （12分）设椭圆 $\dfrac{x^2}{4} + \dfrac{y^2}{3} = 1$，$P$ 为不在椭圆上的点，$l$ 为 $P$ 的极线。
\begin{enumerate}[label=(\arabic*)]
\item 若 $P(0, 0)$，求 $l$ 的方程并解释几何意义；
\item 过 $P$ 任作两条直线分别交椭圆于 $A, D$ 和 $B, C$，证明：$AC$ 与 $BD$ 的交点在 $l$ 上；
\item 利用（2）的结论，若 $ABCD$ 为椭圆内接四边形，$P = AC \cap BD$，则 $P$ 的极线经过 $AB \cap CD$ 和 $AD \cap BC$。
\end{enumerate}

\item （12分）已知椭圆 $\dfrac{x^2}{a^2} + \dfrac{y^2}{b^2} = 1$（$a > b > 0$），$A, B$ 为椭圆上关于原点对称的两点，$P$ 为椭圆上异于 $A, B$ 的动点。
\begin{enumerate}[label=(\arabic*)]
\item 证明：$k_{PA} \cdot k_{PB} = -\dfrac{b^2}{a^2}$；
\item 若 $Q$ 为 $PA$ 的中点，$R$ 为 $PB$ 的中点，求 $QR$ 中点的轨迹；
\item 设 $M$ 为 $AB$ 中点（即原点），证明：$|MP|^2 = \dfrac{|MA|^2 \cdot |PB|^2 + |MB|^2 \cdot |PA|^2}{|PA|^2 + |PB|^2}$。
\end{enumerate}

\item （12分）设 $C_1: \dfrac{x^2}{4} + y^2 = 1$，$C_2: x^2 + \dfrac{y^2}{4} = 1$。
\begin{enumerate}[label=(\arabic*)]
\item 求 $C_1$ 与 $C_2$ 的交点；
\item 设 $P$ 为 $C_1$ 上一点，$l_1$ 为 $P$ 关于 $C_1$ 的极线，$l_2$ 为 $P$ 关于 $C_2$ 的极线，求 $l_1$ 与 $l_2$ 的交点轨迹；
\item 在（2）的条件下，证明：$l_1 \perp l_2$。
\end{enumerate}

\end{enumerate}

\end{document}
